\section{Whole-plane walls, canonical fans, and finite local orders}\label{sec:walls}

This section establishes the local geometric comparison needed when the exterior of an observation region is resampled.  The relevant assertion is uniform over countably many conditional exterior samples: near a buffered compact set there are only finitely many possible local exploration orders, with depth arrays recorded up to their local chart shifts.  Finiteness for each exterior separately would not suffice.  The accounting between these shifts is deferred to Section~\ref{sec:kernels}.

The proof has three geometric stages. We first construct two alternating
wall systems from one scalar GFF and prove that their local arcs are unchanged
under conditional resampling of the exterior. Their half-layer explorations
then give a graph of touching walls with shrinking complementary cells.
Finite shields in this graph bound all relative local depths. Finally, a
local uniqueness argument for directed fan paths upgrades the finite list
of caller geometries to a finite joint order for every ordinary target.
Absolute depth shifts and integer gluing residues remain for the next section.

\subsection{The scalar field and its two wall systems}

Let $\lambda$ be half the height gap for the scalar Dirichlet GFF; with logarithmic covariance coefficient one, $\lambda=\pi/2$.  In this subsection we retain physical height units.  Later integer heights mean multiples of $\lambda$.  Let $J$ be the ordinary scalar whole-plane GFF modulo $4\lambda$, with uniform additive phase, independent of the Liouville surface.  Pullbacks of $J$ have no logarithmic Jacobian term.  Orient a level-$c$ segment so that its left and right data are $c-\lambda$ and $c+\lambda$; reversal together with $J\mapsto-J$ exchanges the choices.

The continuum loop framework comes from the loop-soup construction of
Sheffield--Werner and the sphere nesting of Kemppainen--Werner
\cite{SWcle,KWnested}.  Its GFF interpretation rests on the level-line and
local-set coupling of Schramm--Sheffield \cite{SchrammSheffieldContour}.
The signed explorations and target consistency used here are those of
Miller--Sheffield--Werner and Wang--Wu \cite{MSWTrunk,WangWu}; the alternating
half-height layers follow Aru--Sep\'ulveda--Werner
\cite[Section~6.1 and Remark~19]{ASWlocal}.
Our additional task is local determination under changes of the exterior,
including the finite orders and exact incidences used by the observation
kernels.  The normalization in \cite{ASWlocal} becomes ours after multiplying
its field by $\sqrt{2\pi}$, changing its height parameter $\lambda=\sqrt{\pi/8}$ to
our $\lambda=\pi/2$.

We use the following established disk results.  In a simply connected constant-boundary Dirichlet domain, the two-valued set at offsets $(-2\lambda,2\lambda)$ is the signed nonnested $\CLE_4$: conditional on the loops its signs are independent fair signs, and its fillings are independent Dirichlet fields with the indicated constants.  More generally, the two-valued bounded-type thin local set at $(-a,b)$ exists and is uniquely determined by the field when $a,b>0$ and $a+b\ge2\lambda$ \cite[Proposition~2]{ASWlocal}.  We apply this uniqueness only after proving locality, bounded conditional mean, thinness on compact sets, and the required connection to the boundary.  The piecewise-constant force-point level-line coupling, determination, and target-independence are those of \cite[Theorems~1.1.1--1.1.3 and~1.1.7]{WangWu}.

\begin{proposition}[Whole-plane walls and local determination]\label{wall:local-determination}
There are two signed nested wall systems $O,E$, coupled with $J$, with the following properties.
\begin{enumerate}[label=(\roman*)]
\item Each has the marginal law of a whole-plane nested $\CLE_4$.  Domain means for $O$ are even multiples of $\lambda$ and its wall levels are odd multiples; $E$ has the shifted parity.
\item Under counting Palm measure at a loop, its two side disks carry independent Dirichlet fields with constant side data differing by $2\lambda$, modulo the common period.  The nesting on either side is the deterministic level-CLE construction of that side field.
\item Both systems are determined by $J$.  Under conditionally independent resampling outside an open patch with its field fixed, their arc portions within the patch, together with directions and local data, agree almost surely.  This statement does not fix the outside grouping of separated arcs.
\item The two systems do not cross.  If a wall has a proper open arc in one side of another wall and touches that wall, its level equals the actual mean on that side in a common real lift.  For two conditionally independent copies of the same system given $J$, unequal loops cannot meet.
\end{enumerate}
All assertions about conditional resampling hold simultaneously for countably many sampled conditional extensions and countably many buffered observations.
\end{proposition}

We prove the proposition in several steps.  In particular, a Palm description alone will not be used as a determination theorem.

\paragraph{Signed reconstruction.}
Start with the nested sphere $\CLE_4$ of \cite{KWnested}.  Give loops independent fair orientations, specifying increments $\pm2\lambda$, and choose a fair phase for even domain labels modulo $4\lambda$; their phases alternate across each loop.  Macroscopic local finiteness permits comparison of two loop labels through a finite nesting chain.  Real lifts along that chain retain the signed increments.

Inside any side disk, finite-layer piecewise constant labels are conditional expectations of the Dirichlet field plus its boundary constant.  The residual Dirichlet fillings have variance tending to zero on compactly supported test functions, by mesh shrinkage and Green-function domination.  Martingale convergence reconstructs the field from the signed hierarchy alone.  These reconstructions agree on descendant disks.  Exhausting the plane by increasing side disks defines a distribution $J$ there.

We check compatibility also between overlapping side disks neither of which contains the other.  Suppose their disjoint boundaries each lie in the other's exterior-looking side.  In either disk, the other boundary occurs at a finite generation.  The generation interiors in their overlap are side disks in both reconstructions, so the fields agree there.  The remaining finite-layer carpet cannot support their difference.  Indeed, on every interior compact set, the $\eps$-neighborhood of such a carpet has volume $O(\eps^d)$ for some random $d>0$, almost surely along dyadic scales.  To obtain this, use the exponential tail for the CLE log-conformal-radius decrement \cite{SSWradii}, iterate a fixed finite number of times, apply Koebe distortion, integrate over interior points, and use dyadic bounds.  Both fields have local negative H\"older--Besov regularity of every negative order.  Cover the carpet by $O(\eps^{-2+d})$ mesh boxes and use smooth cutoffs of height one in those boxes.  The scaled GFF test bounds are $O(\eps^{2-o(1)})$, so their total tends to zero.  Thus no discrepancy is supported on the carpet.  The same argument applies to individual loop arcs.  Sums of independent zero-Dirichlet fields on disjoint interiors, extended by zero, obey the same local estimates by Gaussian covariance domination.  In particular, cutting creates no additional distribution on the trace.  All statements across a randomly observed wall transfer from its individual side-Palm law.  Under loop Palm, the two side nestings are independent conditional on the loop, which gives the asserted joint side-field prescription.

The marginal reconstructed field is the required whole-plane GFF modulo $4\lambda$.  Take nesting stops from infinity toward zero at increasing log-radius thresholds, with fresh disk interiors and boundary phase independent of the interior Dirichlet field.  Their interiors exhaust the plane by nesting and local finiteness.  Dirichlet covariances converge on mean-neutral tests to whole-plane covariance.  The variance of a fixed average tends to infinity, while its covariance with neutral tests remains bounded, by the disk Green formula and interior distortion.  Hence its phase modulo $4\lambda$ tends to uniformity independently of the neutral modes.  The compatibility just proved gives the distributional interpretation in a chart about infinity.  This constructs $(J,O)$.  Shift by $\lambda$ and disintegrate to construct $E$ conditionally independently of $O$ given $J$.  We may likewise take conditionally independent copies of either kernel before proving determination.

\begin{lemma}[Arc Palm prescription and localization]\label{wall:arc-palm}
Choose coordinate balls from a countable covering family at generic radii.  Mark the closure $a$ of a proper component of the intersection of a Palm loop with an open ball, with distinct endpoints on the ball boundary, using component counting. Its endpoints are the first boundary hits from the component interior; later returns of the complementary arm belong to other components and are allowed.  Restrictions such as meeting a strict inner collar or having a positive minimum diameter are allowed when measurable from the arc.  Conditional on $a$ and its direction, the remaining arm is chordal $\SLE_4$ in the sphere slit along $a$.  Conditional also on its sign and side data, $J$ on that slit complement is the Gaussian field with the two corresponding constant boundary data.  The assertion persists after extending the observed arc by disjoint nonterminal prefixes.

The conditional intensity kernels of either full loops or these observed arcs, restricted to objects contained with closure margin in an open collar, depend only on the field in that collar.  They do not specify an outside completion or restrictions depending on one.
\end{lemma}

\begin{proof}
We apply the loop-measure identities directly at $\kappa=4$.
Ang--Cai--Sun--Wu identify the counting intensity of whole-plane nested
$\CLE_\kappa$ with a positive constant times Zhan's SLE loop measure for
every $\kappa\in(8/3,8)$ \cite[Theorem~1.1]{ACSWloops}.
Give each loop a fair direction and write $M_\eta$ for its ordinary
$3/2$-dimensional Minkowski-content measure in the coordinate chart.
This geometric measure is independent of any Liouville field.
Zhan's two-mark identity \cite[Theorem~4.2(i), equation~(4.7)]{ZhanLoops}
says that marking with $M_\eta(dz)M_\eta(dw)$ gives, conditional on the
distinct marks, a two-sided whole-plane $\SLE_4$ loop.
Its complementary arm, conditional on either arm, is chordal $\SLE_4$ in
the slit complement \cite[Section~2.2]{ZhanLoops}.

Fix one of the coordinate balls $B$.
Retain ordered marks $z,w\in B$ only when the directed arm $\alpha$ from
$z$ to $w$ lies in $B$.
Conditional on $\alpha$, run the other arm from $w$ to its first exit
$p$ from $B$.  Reverse the remaining chordal curve and run it from $z$
to its first exit $q$ from $B$.  Retain only successful exits before
closure with $p\ne q$.
These are successive stopping operations; the chordal Markov property and
reversibility leave the untouched remainder chordal in the complement of
the observed arc.  That arc is exactly the closure $a$ of the component of
$\eta\cap B$ containing $\alpha$.  The remainder may return to $B$;
no restriction is imposed on such later visits.

For each directed component $a$, the eligible ordered mark pairs are
precisely those ordered along $a$.  Since $M_\eta$ is diffuse, their mass
is $M_\eta(a)^2/2$.
The content $M_\eta(a)$ is a positive finite function of the observed arc
alone.  Weight the marked experiment by $2/M_\eta(a)^2$ and forget the
marks.  This cancels the mark mass and gives component counting, up to the
fixed normalization of loop intensity.  The weight does not change the
conditional remainder law, whose kernel depends only on $a$ and its
direction.  Arc-measurable restrictions likewise leave this kernel
unchanged.

The arc intensity is sigma-finite without a limiting argument.
Indeed, the marginal two-mark density in Zhan's equation~(4.7) is a
constant times $|z-w|^{-1}\,d^2z\,d^2w$ at $\kappa=4$, which has finite
mass on $B\times B$.
Consequently the component-counting intensity restricted to
$M_\eta(a)\ge 1/m$ is finite for each positive integer $m$.
These restrictions exhaust the proper components under consideration.
Further disjoint stopped prefixes have the same resampling rule by the
chordal Markov property and reversibility.

Integrating the full-loop field prescription over that chordal remainder is exactly the ordinary level-line coupling in the slit disk, with its height-gap boundary data.  This proves the Gaussian rule, including observed prefix extensions.  The trace carries no extra distribution by the preceding thinness argument.

Finally, given an observed constraint and its data, the conditional field outside a containing collar, given its field inside, is the ordinary whole-plane conditional field, by the Gaussian prescription and the Dirichlet Markov property.  Bayes disintegration localizes its intensity kernel.  The global period is one common constant: after fixing a collar lift, the extensions are ordinary absolute-height extensions.  Use countably many collars to choose simultaneous versions.
\end{proof}

\begin{lemma}[Product observations and their real lifts]\label{wall:product-gaussian}
Take conditionally independent observation kernels given $J$: a full loop $L$ in one system and a proper initial arc $\sigma$ in another system or copy, with $\sigma$ inside a side disk $D$ of $L$.  Retain directions and labels.  Fix the real lift by the $D$-side trace at $L$ and retain the relative lift at $\sigma$.  Conditional on the two observations, the field on $D\setminus\sigma$ is the Dirichlet GFF plus the harmonic function with those actual boundary values on the two boundary components.  The same prescription holds after prefixes extending $\sigma$, as long as they remain disjoint from $L$.
\end{lemma}

\begin{proof}
Initially put the constraints in disjoint collars with gaps, the arc collar inside $D$.  Their product shape marginal is sigma-finite.  To check the potentially nontrivial integrability, under the first-loop Palm law the field class on the compact arc collar has bounded density relative to its ordinary class, with bound depending on $L$ and hence truncatable on shape windows.  On neutral collar tests the centered Dirichlet covariance is dominated by whole-plane covariance.  Its difference kernel is smooth, and after cutoff on a collar enlargement the relative perturbation is trace class, by arbitrarily high Sobolev smoothing.  The smaller covariance also bounds below a positive multiple of the larger, by Dirichlet-norm cutoff and neutral-test duality.  Finite-dimensional Gaussian determinant comparison, followed by the trace-class limit, gives a bounded centered density.  Conditional on neutral modes, the local average phase has bounded density too: its unwrapped conditional variance has a positive lower bound from a Cameron--Martin shift constant on the observation neighborhood and supported in $D$.  The side-mean shift is constant there.  Combine this bound with Lemma~\ref{wall:arc-palm} and partition into collar and shape windows.  The same argument works for two disjoint full-loop observations when the second is in a constant-data disk.  Prefix extensions need only be disintegrated after the initial finite restriction.

Off the constraints, kernel localization and conditional independence give the ordinary Dirichlet spatial Markov prescription in every strictly interior test domain, conditional on its outside and the observations.  Use a countable exhausting domain library.  Relative lifts are read in collars outside each test.  Exhaustion implies that the conditional field is a mixture of a zero-Dirichlet field and a harmonic function: the independent Dirichlet projections have convergent covariances, and distributional tightness plus interior harmonic bounds make the harmonic remainders tight on compact sets.  The argument also applies jointly to complementary domains.

Identify that harmonic function by traces.  Near $L$ use its side-disk round chart; near $\sigma$ use the round chart of its full slit complement.  Angular distributional traces approach the specified constants or step data almost surely along geometric radii.  For each individual Palm law this follows from the disk Green formula in negative Sobolev spaces, and marginal absolute continuity transfers it to the joint law.  These traces identify absolute and relative lifts from collar data.  A zero-Dirichlet field in the doubly cut domain has zero such traces by covariance domination.  Harmonic measure from a fixed interior compact set has uniformly smooth converging angular kernels in the disjoint conformal collars.  Therefore the harmonic remainder is the bounded harmonic extension of precisely those traces.

One may avoid unsmoothed field restrictions by averaging over geometrically shrinking sleeves.  After subtracting the specified harmonic extension, every angular Fourier mode has zero limit at each circle.  In a circular collar it is a combination of growing and decaying powers, or of a constant and logarithm for the zeroth mode.  Vanishing at the boundary and bounded coefficients on an interior circle force harmonic extension with zero boundary value there.  The exponential mode bounds make this argument simultaneous for the entire distribution.  Apply it at both boundary components.

Repeat with extended-arc observations to obtain the prefix rule.  Initial traces keep lifts consistent.  Rational intrinsic prefix times suffice on disjointness, followed by localization and continuity of conditional means and Dirichlet projections before touching.  Restrictions through a tested time use only the observed prefix, never its future.  The field prescription also holds through the newly drawn trace on compact subsets of the previously open domain: its neighborhood-volume bound transfers from the single-system loop law, and the fields on either side, extended by zero, have arbitrarily small negative regularity by covariance domination and bounded added means.  Thus there is no distribution supported on the trace.

Throughout, all heights belong to one representative of the global real distribution modulo $4\lambda$, fixed by the trace on $L$.  There is no winding freedom in lifting a scalar field.  Two observed arcs of the same oriented loop have consistent actual side data by their individual full-loop law.
\end{proof}

\begin{lemma}[Stopped boundary avoidance]\label{wall:noncontact}
In Lemma~\ref{wall:product-gaussian}, let the continuing arc have level $c$, and let $d$ be the datum on the $D$ side of $L$.  If $d\ge c+\lambda$ or $d\le c-\lambda$, the continuation cannot reach $L$ from inside $D$.
\end{lemma}

\begin{proof}
Suppose it reaches $L$, and follow it to first contact.  In annular coordinates it is a simple slit from its initial tip with open part interior, avoiding the initial arc and the other boundary until that contact.  Cut away a band elsewhere joining the boundary circles, with margin from this slit, to obtain a simply connected domain $T$ agreeing with the annulus near the protected range.  Crosscut topology supplies this cut.  Choose an artificial target away from the protected range and finite piecewise constant data matching the true data there, placing the prospective hit on the right repelling flank if $d\ge c+\lambda$ and on the left if $d\le c-\lambda$.  All right data can be at least $c+\lambda$, and all left data at most $c-\lambda$.  The path is admissible as a Loewner slit aimed at that target through its protected interior prefixes.  A countable library of cuts and neighborhoods in the annular chart covers all tests with strict margins.

We justify changing the Gaussian domain under one fixed positive tilt.  Put two separated nested comparison collars between the protected range and the artificial changes, allowing them to meet shared true boundaries.  For a fixed protected prefix $a$, compare four prescriptions: the annular and $T$ base fields, and the corresponding fields with $a$ and its level jumps inserted.  Their restrictions to the two collars are equivalent to a common product local Gaussian class.  On enlarged disjoint collars the Dirichlet Markov pieces agree, and the remaining harmonic differences have finite Dirichlet cutoff norm.  At a shared true boundary they vanish and obey the same boundary energy estimate in the analytic collar chart.  Tips and jump discontinuities are common or separated from these collars.  The two base laws are also equivalent through the protected interior; the two inserted laws are equivalent there with the common slit.

Let $R$ be the $T$-to-annulus base density and $R_a$ the analogous inserted density.  By the Markov property both are functions of the outer collar, even when the protected field is included.  On the two collars, insertion-to-base densities in each fixed domain use only the inner collar.  Hence $R_a/R$ is measurable both from the outer and the inner collar.
Under a product measure, a random variable measurable from both coordinates
is almost surely constant; equivalence gives the same conclusion here.
Thus, for almost every observed prefix $a$,
\[
 R_a=c(a)R
\]
on the collar restrictions, with $0<c(a)<\infty$. The scalar $c(a)$ may
depend on the prefix, but the base density $R$ does not. These are
conditional density identities under each valid marginal prefix law;
countably many stopped-prefix tests suffice. They do not require one
exceptional-set convention for every deterministic slit.  Tilt the whole joint field-and-path law once by $R$.  Conditional on every tested prefix, normalization cancels the scalar $c(a)$,
so its protected field has exactly the $T$-slit prescription.  Extend to the rest of $T$ by the base conditional kernel, which is unchanged by an insertion given the protected field.  No insertion-to-base density is evaluated on the growing trace.

Stop inside buffered protection, at compact capacity localizations, and before incompatible boundary approaches.  The conditional mean martingales have bracket equal to the Dirichlet Green decrement.  The force-point level-line characterization
\cite[Proposition~2.1.7]{WangWu} therefore identifies the stopped continuation
with the ordinary level line in $T$.  Its cumulative force weights on both chosen sides are nonnegative, including the boundary-avoidance threshold zero; it avoids the interior of the indicated repelling arc.  Its Loewner tip and capacity are continuous before contact, by simplicity and protected domain agreement.

For completeness, all choices of $T$, target, data, collars and protection can be made countably conditional on the initial two shape observations and labels.  Restrictions up to the stops follow first at rational intrinsic times with disjoint constraints, then by continuity or martingale approximation.  One can alternatively continue at a stop with the ordinary Gaussian level coupling in the remaining $T$ target component.  Its prescribed field and admissible cumulative weights give the same mean and Green-bracket martingales for the concatenation; hidden stopping information may then be forgotten.  Thus probabilities of protected approaches are bounded by those of the ordinary $T$ line under the same initial positive tilt, rather than a sequence of densities depending on approach distance.  Exhaust the compact repelling subarcs and then the countable tests, and remove the positive tilt.  First contact is impossible.
\end{proof}

\paragraph{Noncrossing and determination.}
Arc tests inside either side see every proper interior part of an unequal intersecting loop.  Follow the remainder from such a test and apply Lemma~\ref{wall:noncontact}.  For systems of the same parity the continuing level differs by at least $\lambda$ from either side mean, so unequal meeting loops are impossible.  Equal physical loops must have identical actual data by side traces.  They cannot belong to opposite parity systems.  For opposite parity systems both sides cannot permit contact, and a touching loop wholly on one side must have level equal to that side mean.  This proves the noncrossing and label assertions of Proposition~\ref{wall:local-determination}.

Compare now two conditionally independent copies of $O$.  Observe a loop-side disk $D$ of one, with mean $c_D$, strictly enclosed in a parent disk $P$ of the other, with mean $c_P$.  Joint counting is localized as above, with relative lifts retained.  Initially the field on $D$ has its Gaussian side-disk law.  Reveal the direct child carpet in $P$.  Either one child contains $D$, equality allowed, or the child interiors lying in $D$ fill it up to their carpet.  Call the latter alternative \emph{penetration}.  There is no crossing alternative.  In the first case the field on $D$ keeps its Gaussian law; in the second the child interiors carry independent Dirichlet fields with constants $c_P\pm2\lambda$.

Here the conditional assertions follow from the same argument as Lemma~\ref{wall:product-gaussian}, not from independence after observing two systems.  The conditional kernel of the parent with its child layer is local off compact subdomains inside each fresh interior, by its single-system CLE/GFF rule and Bayes.  Compact Markov exhaustion yields a Gaussian-harmonic mixture.  Unchanged individual side traces identify its harmonic part in an uncut $D$, or jointly in the child interiors on penetration.  The carpet supports no residual distribution by the volume and regularity estimates.  This remains valid if child and test-disk boundaries touch, since traces from within each child and strictly interior Markov tests still suffice.

On penetration both child signs occur within $D$ when their boundaries are disjoint from $\partial D$.  Some child boundary lies inside $D$, and each neighborhood from its other side meets infinitely many further child interiors, by disjointness, filling of Lebesgue points, and local finiteness.  Almost surely, both independent signs occur in every rational ball meeting infinitely many children.  If $c_P\ne c_D$, the two possible child values lie on one weak side of $c_D$, and one lies strictly to that side on penetration.  The conditional mean, tested against nonnegative interior functions, then contradicts its initial value $c_D$.  If $c_P=c_D$, plug a fresh canonical two-valued construction into the unchanged nonpenetrating alternative.  The resulting coupling inside $D$ is a two-valued BTLS at offsets $\pm2\lambda$: the just-proved prescription gives locality, bounded means and independent Dirichlet fillings, and the carpet is thin.  Its connection to the boundary follows by detouring a path to $\partial D$ along rims of the disjoint Jordan holes.  The holes have a null sequence of closure diameters; replacement arcs therefore shrink at accumulation points, and each fixed rim is Jordan.  Stop upon reaching $\partial D$, including at a touching rim.  This also verifies the required connectivity in the analogous carpet couplings below.  BTLS uniqueness identifies this carpet with the tested system's own children in $D$.

Every bounded-side $D$ is enclosed by some loop of the other system.  Either $D$ itself is eventually met, with its deterministic children shared, or descending mesh shrinkage forces penetration at a strictly enclosing parent.  The preceding argument again shares all its children.  Symmetry and nesting give equality of the entire two systems.  Thus each conditional kernel is a point mass: $O$, and likewise $E$, is a measurable function of $J$.

Finally, take an arc-Palm observation with collar inside an open patch.  Exterior resampling conditional on the patch preserves the joint law of the arc and field by Lemma~\ref{wall:arc-palm}.  Determination puts that same arc on the resampled system.  Countable ball-and-arc covers and symmetry identify all local portions, including their directions and side data.  This proves Proposition~\ref{wall:local-determination}.  Its conclusion concerns typical conditional extensions, which is exactly what will be used below; it asserts no rule for arbitrary exceptional exterior modifications.
\subsection{Alternating canonical fans}

From now on scalar heights are in units of $\lambda$.  A disk always specifies a side of its boundary.  Neither quantum length nor exploration speed determines the half-layer numbering.  For a Jordan Dirichlet disk $C$ with constant value $c$, write $\cA(C)$ for its two-valued set at offsets $(-1,1)$, including $\partial C$ in the closed set.  Its holes are called \emph{half-cells}.

Two composition identities will be used.  Filling exactly the $c+t$ half-cells, $t\in\{-1,1\}$, with ordinary $(-2,2)$ sets gives the two-valued set at $c-t,c+3t$.  Iterating half-layers until the labels reach $c\pm2$, renewing the procedure in cells returning to $c$, gives the ordinary $(-2,2)$ set.  Finite versions follow from local-set composition and BTLS uniqueness: the fillings are Dirichlet, the means remain bounded, and the required two boundary values are those stated.  For the second identity, one can stop after one pair of layers and insert an ordinary $(-2,2)$ construction in the returns \cite[Section~6.1 and Remark~19]{ASWlocal}.  These almost-sure field identities remain valid under absolutely continuous biases from previously observed containing geometry.

\begin{lemma}[The canonical fan]\label{wall:fan}
In a constant-data Jordan disk $C$, take a source $x$ and a dense countable collection of ordinary targets independent of the Dirichlet field.  Draw the level-$c$ branches, of chordal law $\SLE_4(-1;-1)$.  Their interior excursions are identical or disjoint full simple crosscuts, with only finitely many distinct excursions of diameter above any positive threshold in closed disk coordinates.  Their closure, together with $\partial C$, is $\cA(C)$, independently of the source and the chosen dense target test.

Branches from the same source agree until their targets separate, and any two branches meeting a point agree up to that point.  Their contacts on either boundary flank progress monotonically toward their target.  Every open boundary interval contains excursion endpoints of a crosscut lying arbitrarily close to that interval.  For a specified boundary prime point, the probability that it lies on any half-cell closure is zero.  All statements hold simultaneously for countably many ordinary sources or additional target tests.
\end{lemma}

\begin{proof}
Target-independence in the level-line coupling gives branching up to separation; here the total force weight is $-2=4-6$.  After separation, branches remain in the distinct target components and do not revisit their shared past.  The open original boundary intervals toward the targets are disjoint.  Simple-curve separation then gives the meeting and monotone-contact assertions.  We use ordinary targets only, so no rule is needed at a target lying on another path.

Along a trunk use the full MSW exploration, tracing the encountered $c\pm2$ CLE loops.  When a target is cut into a negative domain, continue from its pinch by the corresponding exploration there.  Proposition~\ref{cut:branching-fields} gives constant-$c$ fresh domains for this full signed exploration.  The trunk alone need not have constant data in every domain it cuts.  Splitting toward two targets and then exploring their disjoint fresh components gives the same branching coupling, with targets on transmitted open arcs.

To prove macroscopic excursion finiteness, equip $C$ in conformal coordinates with an independent auxiliary critical quantum disk.  Initially take a typical length source and target and draw one full MSW branch.  In each negative child draw another such branch from its pinch toward an independent uniform length target; repeat only in negatives.  The one-branch marked rule, negative moment contraction and composed diameter estimates in Lemma~\ref{cut:tiling} imply summable large-domain bounds after a fixed parent lower-proxy cutoff, and uniform diameter shrinkage down these negative generations.

An original typical target is reached by following helper prefixes and switching to its negative child at separation.  The subsequent continuation lies in that child's closure.  There are no corner ambiguities for countably many typical length tests.  Original-boundary contacts are linked outer records; running-record sets and jump endpoints have zero length by Proposition~\ref{cut:critical-cutting}.  Open original arcs transmit isometrically, and loss and aperture boundaries are the prescribed jump boundaries.  Thus a target entering a negative retains an open original boundary interval.  It cannot enter a positive disk, since positive loops do not touch the true boundary of the current negative.  One may equivalently route a fresh independent target after conditioning on the auxiliary surface and helpers, or sample in surviving intervals.

Only finitely many helper cells matter above a prescribed diameter.  Once the first sufficiently small descendant is entered, its entire remaining continuation is small; before then there are finitely many cells by generation and sibling shrinkage and containment.  Every macroscopic excursion therefore contains an interior subarc of definite positive diameter on one of finitely many helper trunks: truncate its small terminal tail and use the finite number of preceding switches.  Distinct excursions have disjoint interiors, so uniform continuity and injectivity on the finite helper list give finiteness.  The same construction shows that each countably tested interior passage belongs to a full excursion, with no extra interior limiting trace.

Individual paths have no boundary arcs; for example their positive jumps are dense and have interior tips by the exact stack rule.  Their hits on either open boundary flank have harmonic measure zero.  To see that every prescribed point other than the endpoints has zero hitting probability, first use the auxiliary length experiment: linked original-arc real-time heights lie in closed descending record sets, of length zero.  Conditional on endpoints and surface the chordal law is conformally invariant; for fixed endpoints its one-point hitting probability is constant on either open flank.  The diffuse full-support length measure forces that constant to vanish.  Boundary contacts occur arbitrarily near the target on both flanks before terminal time: both coordinate records take positive values tending to zero at times before termination, by compact positivity, the c\`adl\`ag rule, and first entries below vanishing thresholds.  Their linkage to the original arcs is exact.  A dense target in a given open boundary interval and late gaps in its contact-time set consequently provide arbitrarily small excursions with both endpoints there.  The contact-time set has empty interior, and simplicity makes those late excursions small.

These conclusions transfer to fixed $x$ by conformal invariance, or by rotating the independent auxiliary chart.  An additional ordinary target contributes no new excursion: an excursion on its path remains on paths toward an open target interval, and a dense test already meets that interval by target-independence.  Thus dense tests can be changed, and the conclusions transport to closed Jordan disks.

We identify the complementary components and the field.  A null sequence of disjoint full crosscuts, with disjoint open portions, has Jordan complementary components, each touching the outer boundary.  For a point outside their closure, remove the shadows cut off from it by each crosscut.  These shadows are nested or disjoint.  Every positive-sized shadow is behind a maximal one: nested enlargement uses only finitely many non-small crosscuts, whereas small crosscuts away from the point have small shadows in Jordan coordinates.  Replace the corresponding disjoint open intervals of the outer circle by their maximal crosscuts in cyclic order.  Their diameters tend to zero, so this is a simple closed boundary with the expected interior.  Coincident endpoints merely make successive arcs abut; if both complementary intervals between two endpoints are replaced, one obtains the two-sided Jordan lens.  No other crosscuts accumulate locally on the open interior of a fixed one, so it has a complementary component on each side.  The uncut case is included.

Finite branch unions are thin GFF local sets with conditional mean in $[c-1,c+1]$: every continuation occurs only in its remaining component, with untouched data $c$ and matching slit data.  Martingale convergence and Dirichlet covariance exhaustion give the prescription off the closed union.  It is thin with bounded mean.  On an interior compact set only finitely many excursions occur, lying on countably many chordal SLE traces of dimension strictly below two; alternatively, bounded finite-stage means converge against every compact test by dominated convergence off the zero-area closed union, while Green covariances decrease to the component covariances by boundary regularity.  All interior boundary arcs facing a given component carry the same value $c+1$ or $c-1$.  Two supporting paths after divergence cannot both face that component, and on one path a connected complementary region lies on one side.

A component's contacts with $\partial C$ lie among the boundary points hit by the countable paths, with at most one additional point.  Otherwise a tested path with endpoints alternating with two unhit points would separate nearby interior points of the component.  These contacts are harmonic-null both in $C$ and within the component.  Brownian exit and local agreement on the crosscuts therefore identify its bounded harmonic mean with the indicated constant.  Boundary connection and BTLS uniqueness yield $\cA(C)$.  Uniqueness makes its interior crosscut set source-independent; adjacent signs fix crosscut directions.  Finally rotational invariance and harmonic-null contacts, summed over the cells, show that a specified prime point lies on no half-cell closure.  This applies conditionally to countably many points selected before sampling a fresh half-layer.
\end{proof}

\paragraph{Comparing the two parities.}
Fix an $O$ side disk $D$ with mean $c$ and a real lift, viewed from an outside point.  There is an $E$ parent $P$ containing $D$, with possible boundary contact, at whose next child layer penetration occurs and no child encloses $D$.  Follow a nesting chain separating a typical interior point from a sufficiently exterior point.  Large side disks enclose $D$, and shrinkage with noncrossing gives the first penetrating parent.  Rational test points off the walls suffice.  An enclosing child, if present, is unique.

Opposite-parity walls cannot share an open subarc.  Along such an arc their side traces would have to coincide but would have different parity.  This trace comparison remains valid in the random charts: transitions between round side charts extend analytically across a strict common arc by reflection.  More explicitly, average in a sleeve of normal distances $\eps$ to $2\eps$ from a round boundary.  For a smooth tangentially bounded test with density $O(\eps^{-1})$ its zero-Dirichlet variance is $O(\eps)$ by the Green bound.  In fixed tangential and rescaled transverse coordinates, summing variances in a sufficiently negative Sobolev norm gives uniform convergence to zero, almost surely on dyadic sleeves, over tests with bounded sufficiently many derivatives.  This covers both analytic charts without selecting an exceptional random approach.  Palm transfer gives the asserted trace equality.

The penetration argument above also applies with touching $P,D$.  For its disintegration use probability selections of $P$: the first loop inward across a specified log-radius threshold along a chain separating two specified points, followed by a fixed number of generations.  Varying this countable data covers all candidate parents.  Selection precedes revealing the interior, so the marked interior and child layer have their fresh disk rule.  Their kernels are local off strictly interior test domains, whether or not the exterior selection is local.  Product with the other system's finite-window $D$ kernel is finite, and single-system assertions transfer by marginal absolute continuity.  Thus initially $J_D$ has its stated Gaussian law; after revealing the children it is unchanged if a child encloses it, or has the thin bounded local-set prescription with values $c_P\pm2$ on penetration.  Traces inside the child domains and compact Markov exhaustion suffice even at touching boundaries.

If $|c_P-c|\ge3$, conditional expectation excludes penetration.  Otherwise $c_P=c+t$, $t=\pm1$.  Insert an ordinary construction on the nonpenetrating alternative.  BTLS uniqueness identifies the penetrated carpet with the canonical two-valued set at $c-t,c+3t$.  These identities are obtained before restricting to penetration and hence persist on that event.  Exhausting the parent selections proves them for every base observation in use.

\begin{definition}[Canonical paired sweep]\label{wall:paired-sweep}
Let $G_D$ be the closed disk $\overline D$ minus its open direct $O$ child interiors.  Within $D$, and subsequently within each return disk $C$ of value $c$, perform the following two half-steps.
\begin{enumerate}
\item Draw $\cA(C)$.  Its $c-t$ holes are complete $E$ islands of the penetrating layer, all touching $\partial C$; within the $c+t$ holes that layer continues by ordinary CLE.  The entire first fan consists of island-wall arcs, since each crosscut faces both signs.  It lies in $G_D$ by the ordinary two-layer identity.  In either kind of half-cell $H$, with value $h=c\pm1$, the available $E$ nesting is the ordinary level CLE of its field.
\item Draw $\cA(H)$.  The cells with value $c+2(h-c)$ are terminal original $O$ children.  The cells with value $c$ are returns $C'$, with ordinary remaining $O$ construction.  Every new interior crosscut lies on a terminal wall.  Such walls touch $\partial H$ away from $\partial C$, because the original children of the ordinary $O$ construction in $C$ are compact inside $C$.  They are precisely the walls hit by the first fan; deeper $O$ disks within returns cannot touch it.
\end{enumerate}
Within $C'$, the stopping layer of ordinary $E$ from $H$ at $h\pm2$ factors through $\cA(H)$ and restores the first-step situation with $t=h-c$.  Its direct disks meeting the open $C'$ lie there, and no preceding ancestor wall enters it.  Iterate.
\end{definition}

The construction consists of countable compositions of Dirichlet-field identities, so the preceding initial biases are harmless.  Pure $J$-geometry assertions may be proved under the single $O$ base Palm law, where each canonical half-layer is sampled from its fresh Dirichlet filling.

\begin{proposition}[Returns and finite capture]\label{wall:returns}
For each return $C'$ born in a paired sweep of $C$,
\begin{equation}\label{wall:return-compact}
\overline{C'}\subset C.
\end{equation}
Every return is captured as one untouched negative pocket by finitely many ordinary full-fan branches from $\partial C$, and has empty inherited open target interval.  Conversely a complete negative birth with empty inherited target interval is a stabilized return.  As the finite target lists increase densely, the two-layer walls stabilize exactly on interior compact sets and the uncaptured negative pockets have uniformly vanishing closure diameters.
\end{proposition}

\begin{proof}
Use an independent auxiliary critical disk in $C$ and countably many ordinary full-fan paths.  Suppose $p\in\overline{C'}\cap\partial C$.  Then $p\in\partial H$.  Apart from at most one point for each $H$, these contacts are visited by a tested trunk.  Also exclude all test endpoints and all endpoints, in trunk time, of boundary-contact gaps.  This countable exceptional set is visible from the first fan and is avoided by every next half-cell closure by the specified-point assertion in Lemma~\ref{wall:fan} and conditional Dirichlet filling.  Thus $p$ is a nongap boundary visit at a real time $0<v<T$ on an open original side arc of one tested trunk.

The connected return avoids that trunk and all hit walls and is not inside a positive disk.  It lies in a single negative jump disk of the full branch.  Consequently $p$ has three representations: an original open-arc label, the real tip at $v$, and a point on the loss circle of a negative jump $u$.  Apply the exact stack and polarity rules of Proposition~\ref{cut:critical-cutting}.  The original arc links backward from $v$ by a sustained running record on a real copy of one side $\beta$.  If $u\ne v$, the loss representation is not a second real tip, by the no-two-real-times rule.  It is an open loss label linked forward from $v$, so $u>v$.  The side and real copy are the same, since a tip cannot carry two nontrivial links.  Until $u$ its height stays above the held label.  Earlier equality and simplicity exclude further original-arc record contacts in the between-interval.  Other-side contacts cannot accumulate at an interior point of this original side arc.  Hence $p$ is a gap endpoint, a contradiction.  If $u=v$, the real-time record convention puts $\beta$ on the jumping side and its label at the post-jump value.  Contacts cannot converge from the left, whose heights instead tend to the distinct pre-jump value.  This again makes $p$ an excluded gap endpoint.  This proves \eqref{wall:return-compact}; it is also consistent with the two-valued-set boundary-avoidance statement of \cite[Lemma~6.11]{ASlocal} at $r=\lambda$ and its sign reversal.

On an interior compact set, finitely many branches capture all first-layer crosscuts completely by Lemma~\ref{wall:fan}.  Their hit-loop walls are then drawn.  A terminal wall reaching a smaller compact set but attached farther away is macroscopic; only finitely many exist, and each is eventually drawn by the same argument.  Thus the full two-layer set stabilizes exactly on compacta.  By \eqref{wall:return-compact}, the entire boundary of any return is captured by finitely many branches while its interior is avoided.

The resulting pocket is one actual negative cut, not a union of leaves.  After a branch separates into an earlier negative, it shares the prefix to separation and then explores only that untouched component; each new negative is a single cut with boundary on the full branch or its old outer boundary.  Positive interiors are put aside.  Inductively every untouched nonpositive component after finitely many branches is one such negative.  Capturing the return's full Jordan boundary while avoiding its interior identifies exactly one pocket.  Its birth is shared by all branches through its cut node; branching elsewhere cannot rediscover it.  It has no surviving original open target arc.  Conversely a complete negative with no such arc is never entered again by typical original targets, hence persists as a complementary component and is a return.  Its closure is therefore disjoint from the original boundary.

For uniform shrinkage of remaining negatives, first work away from $\partial C$.  After compact wall stabilization, each return meeting a smaller compact is captured: all but finitely many are small and lie in the stabilized region, while the two successive half-layers have only finitely many cells of each positive diameter.  A sequence of remaining macroscopic pockets in boundary collars of vanishing width would have a connected Hausdorff limit containing a boundary interval.  Choose two unhit points there alternating with the endpoints of a tested path.  That path separates their small neighborhoods, contradicting a later pocket's interior traversal between them.  Positive walls above any diameter are finite and are eventually reached whenever eligible.  This proves the assertions and realizes each sweep by refining finite ordinary branch lists.  No infinite length-conditioning is used at a birth.  On the lattice, a positive-size captured return has no surviving original target slots for all large indices, by compact containment and the upper topological approximation of transmitted occurrences from Section~\ref{sec:cutting}.
\end{proof}

\begin{lemma}[Half-layer mesh under conformal distortion]\label{wall:mesh}
There are constants $A_0,d>0$ for fixed tent conventions such that, in a round prime-disk chart, a boundary tent of width and depth $O(t)$ is met by a new half-cell closure of diameter at least $2^jt$ with probability at most $A_0e^{-dj}$.  If $f:\Delta\to C$ is a Jordan conformal map into a bounded plane chart, then
\begin{equation}\label{wall:image-mesh}
\Prob\{\text{some new half-cell image has diameter greater than }r\}
\le \min\{1,A_0|C|/r^2\}.
\end{equation}
In a bounded chart for a base disk, closure diameters of active canonical domains tend uniformly to zero almost surely.  At each fixed half-step only finitely many cells exceed any prescribed positive diameter, and every original $O$ child is reached at finite depth.
\end{lemma}

\begin{proof}
Use an upper half-plane chart at the tent center.  The constant-boundary chord from $-s$ to $s$ is in the half-layer by source-independence.  For a fixed large $R$, with probability at least $p>0$ it stays, including its endpoints, in $s/R<|z|<sR$.  This follows from prescribed-point avoidance and compactness at unit scale.  A successful chord at an intermediate scale separates the tent, including nearby points used to test closure contact, from the far region.  Close the crosscut outside the domain along the base interval.  Additional boundary contacts are harmless since the chord avoids the enlarged tent ball.

Success is measurable from the field in a slightly enlarged band.  Given the successful chord, its conditional exterior field given the band is unchanged by the Dirichlet Markov prescription; Bayes and path determination localize the successful-path kernel.  In logarithmic coordinates the field is zero-Dirichlet on the two sides of a strip of height $\pi$.  Test $m$ translated bands separated by a large fixed gap $g$.  If $q$ is their joint density relative to the product of their marginal restrictions, then
\begin{equation}\label{wall:band-density}
\int q^2\dd\Prob_{\mathrm{prod}}\le\exp(Cme^{-2g}).
\end{equation}
Indeed the sine modes are independent stationary Ornstein--Uhlenbeck processes of rates $n=1,2,\ldots$.  By the Markov property, dependence between intervals is determined by endpoint covariances.  Normalize each mode's variance to one.  Equal-interval endpoint covariance matrices have uniformly bounded inverses; off-block perturbations have operator norm $O(e^{-ng})$ by geometric row sums and squared Hilbert--Schmidt norm $O(me^{-2ng})$.  Diagonalizing relative to independent blocks, a relative covariance eigenvalue $1+\delta$ contributes $(1-\delta^2)^{-1/2}$ to the density's second moment.  Multiply finite-dimensional determinants, then take the product and $L^2$ limit, summing over $n$, to obtain \eqref{wall:band-density}.  Cauchy--Schwarz bounds the probability all chords fail by
\[
\exp(Cme^{-2g}/2)(1-p)^{m/2}.
\]
Choose $g$ large.  There are order $j$ such bands between the two scales, proving the exponential tent bound, with constants enlarged for finitely many chart-edge scales.

For \eqref{wall:image-mesh}, let $a_Q$ be the conformal energy of $f$ in a fixed enlargement of a Whitney box $Q$ of scale $t$, including finitely many neighbors.  Then $\sum_Qa_Q\le C|C|$, and the image oscillation in the box is at most $C\sqrt{a_Q}$ by holomorphic interior estimates.  At relative scale $j$ call $Q$ dangerous if
\[
\sqrt{a_Q}>c_0r/(j+1)^2.
\]
Exclude the event that a half-cell of prime diameter between $2^jt$ and $2^{j+1}t$, up to fixed factors, meets the tent above such a box.  Include finitely many neighboring tent conventions and the central order-one boxes.  The tent estimate and energy sum bound the union probability by
\[
\frac{C|C|}{r^2}\sum_{j\ge0}(j+1)^4e^{-dj}\le \frac{C'|C|}{r^2}.
\]
Outside that event, compare any two closure points of a half-cell, which touches the boundary, by radial Whitney chains up to boundary distance comparable to its diameter, and a bounded horizontal connection there.  Each used box has an associated tent near a closure point, with bounded multiplicity at each relative scale.  Sum the oscillation bounds and use boundary continuity of $f$.  With $c_0$ sufficiently small the total is at most $r$.

At any fixed area point, the labels in a fresh half-step move by $\pm1$ with equal probability, by its bounded conditional mean.  The active process stops at $c\pm2$.  Its survival probability, and hence expected total active area, decay geometrically with half-layer number.  Conditional on the earlier domains, apply \eqref{wall:image-mesh} to every new half-layer and sum over domains and steps.  Borel--Cantelli, first for rational $r>0$, gives uniform closure-mesh shrinkage.  This proof is under the ordinary Dirichlet base law and transfers to the wall comparisons.  At a fixed stage, component closures form a null diameter sequence.  Every original $O$ child is terminal at finite depth, also because its positive-area open interior must be filled at some finite step.
\end{proof}

\begin{lemma}[Wall graph and finite-arc approximation]\label{wall:graph}
In $G_D$ use as vertices the full loops $\partial D$, the direct $O$ children in $D$, and the $E$ walls entering $D$ and lying wholly in $G_D$.  Join touching vertices.  Then the closed set drawn after $j$ half-layers is exactly the union $S_j$ of vertices at graph distance at most $j$ from $\partial D$.  There is no same-kind adjacency.  Each $S_j$ is compact and has the following local connection property: points tending to one of its points can be joined to it by finitely many actual wall arcs of diameter tending to zero.

Every path in $G_D$ whose endpoints lie on reached walls can be approximated, in path order and in an arbitrarily thin neighborhood, by a path made of finitely many wall arcs.  For a closed path the approximation preserves winding about a smaller protected set.
\end{lemma}

\begin{proof}
An $E$ wall with an open portion in the base outside the children, on their exterior sides and not contained in a child boundary, cannot leave $\overline D$ or enter a child elsewhere, by noncrossing.  A contact visible only from the other side of an $O$ wall is not an edge in this base.  The first islands are precisely the eligible neighbors of $\partial D$: deeper $E$ disks in the penetrating interiors cannot reach it.

A terminal $O$ child in a half-cell $H$ touches the preceding $E$ arcs and avoids the rim of its current return $C$, because the remaining $O$ law there is ordinary.  Conversely an undecided $O$ child touching a newly introduced island lies in the same $C$.  Its contact cannot be on an interior-stage $\partial C$, made of old $O/E$ pieces, without a forbidden same-kind unequal meeting; it cannot hit $\partial D$ either.  Deeper $O$ disks inside returns cannot reach the new islands.  Exchange the kinds at an added terminal $O$ wall.  An unreached eligible $E$ wall touching it does so within the same $H$: the terminal avoids $\partial C$, and the rest of $\partial H$ consists of new $E$ pieces.  Ordinary $E$ nesting in $H$, split by its half-cells, shows that this wall must be direct inside an adjacent return $C'$.  Deeper disks cannot reach the dividing boundary.  Conversely a new $E$ island in $C'$ touches a terminal: it touches $\partial C'$, cannot meet an earlier unequal $E$ wall, and $\overline{C'}\subset C$.  These are exactly breadth-first additions.  Any eligible remaining $E$ wall either touches a reached vertex and is reached next, or persists in active domains forever, contradicting Lemma~\ref{wall:mesh}.  No nonlocal sign choice is used to resolve a tie.

Compactness of $S_j$ follows from attachment and loop local finiteness.  Prove the local connection property inductively.  Small new walls attach to $S_{j-1}$ nearby; the fixed finite collection of larger new walls is controlled by Jordan continuity and compactness.  This gives connections by finitely many actual wall arcs with vanishing diameter.

For path approximation take $j$ large enough to contain the endpoints and to make every active hole tiny.  Each excursion of the proposed path outside $S_j$ lies in one active open hole, never in a terminal interior.  Its entry and exit join near the hole boundary by finite wall-arc paths, by the local connection property and boundary connectedness.  To avoid an infinite concatenation, split the desired path into short pieces in small balls.  On its closed set of hitting times of $S_j$, joinability by finitely many arcs in an enlarged ball is an open equivalence relation.  Successive complementary-interval endpoints are equivalent by the preceding hole argument.  Compactness and interval connectedness give one class, using a finite chain.  Macroscopic passages admit subdivision endpoints on $S_j$ because the holes are tiny.  Concatenate the resulting finite approximations.  Keeping the neighborhood off a protected set preserves winding.
\end{proof}
\subsection{Finite local cell patterns and source pools}

The preceding graph approximation supplies finite wall-arc paths surrounding
small protected regions. We use them as shields: any exploration reaching
the protected region must meet one of their finitely many wall vertices.
This converts possibly unbounded exterior exploration histories into bounded
relative depths. The remaining issue is the order in which fan paths use
the resulting local geometry; that is addressed after the source coupling.


Fix a compact observation set $K$ buffered strictly inside a chart neighborhood $U$.  Compare countably many conditional extensions with the same local $J$.  The distinguished starting $O$ ring separates two marked points outside the local buffers.  Make canonical sweeps on both sides of it, then in reached bases and their return hierarchies.  Each compared sample retains its valid marginal, including the local wall-Palm rules.  All equalities below are almost-sure equalities for these sampled extensions and ordinary path tests.  Small comparison disks contain neither marked point.  Remote-only steps may be arbitrarily numerous; the assertions concern local operations, and sometimes record depth differences only.

\begin{lemma}[Finite shields and relative depths]\label{wall:shield-depth}
Around every point of $K$ one can choose a smaller comparison neighborhood with common local side-gasket geometry, or identify an entire common internal $O$ subtree.  In a side-gasket neighborhood there are finitely many shared wall arcs forming a shield.  If their depths in a compared sample are $d_1,\ldots,d_m$, then for every wall occurrence $w$ in a smaller protected neighborhood,
\begin{equation}\label{wall:shield-formula}
\operatorname{depth}(w)=\min_{1\le i\le m}
\{d_i+\ell_{\mathrm{loc}}(i,w)\}.
\end{equation}
Here $\ell_{\mathrm{loc}}$ counts graph edges along actual wall-arc connections inside the common comparison neighborhood.  The vector of depth differences has finitely many possibilities.  Thus the whole inner array of wall depths has finitely many possibilities up to one common even shift and finite parity data.
\end{lemma}

\begin{proof}
Work first in a small disk $B\Subset U$ about $z$.  If $z$ lies in the compactly contained bounded side of a complete $O$ wall, that wall and side are common by Proposition~\ref{wall:local-determination}.  Its side is explored inward, since the starting ring cannot be within it; ancestor gasket sweeps do not enter it.  Its entire canonical interior hierarchy is common.

Otherwise only finitely many nonlocal $O$ arc pieces can approach $z$.  Pieces crossing a collar are finite by loop local finiteness and Jordan continuity; all other nearby walls are complete within $B$.  Unequal $O$ walls are disjoint.  If $z$ is off the walls, it lies in a base gasket away from each individual child closure.  Infinite nesting at $z$ would have put it inside a complete small wall, so such a base exists.  Its boundary lies outside a smaller neighborhood.  The only child boundaries nearby are known complete loops, namely maximal small loops on the outward side toward $z$.  Nonlocal arcs remain away and cannot separate nearby wall points from $z$; a complete loop compactly inside $B$ has its bounded side inward.  Thus the omitted open holes are known locally.  If $z$ is on a wall $W$, shrink to a Jordan box where only its arc, among non-small $O$ walls, approaches $z$.  On each side use the adjacent base gasket, where $W$ is the base rim or a direct child.  Other nearby holes are complete and local.  One may first handle larger loops in a larger box and require each loop called small to have its bounded closure inside $B$.

The relevant actual mean $c$ is fixed by a common $O$ side trace: by $W$ in the on-wall case, or an exposed nearby child arc in the off-wall gasket case.  Such arcs approach gasket points by filling and disjoint Jordan nesting.  Therefore $J-c$ agrees as a real field, rather than merely after independent choices of representatives modulo the period.  The restricted graph is also common.  An $E$ arc facing the open base belongs according to its side relative to the known holes; noncrossing prevents a remote continuation into a different side.  Read contact-only points by continuity along arcs, since opposite parities share no open subarc.  Zero-cost movement follows actual local arc connections and does not assume an outside grouping of separated pieces.

Construct the shield off walls by starting with a small surrounding path and detouring along hole rims.  Large holes stay away from $z$ and remaining holes are arbitrarily small, so the detours exist arbitrarily close to $z$, remain at positive distance from it, and retain winding by homotopy off $z$.  Jordan continuity controls accumulating gaps.  Lemma~\ref{wall:graph} approximates this path by finitely many wall arcs in one sample, preserving winding about a still smaller neighborhood.  These arcs are shared in every sample.  At a wall $W$, use a shield arch from $W$ to $W$ through the chosen side and away from $z$, obtained from a semicircle in Jordan coordinates with the same detours.  Close it along $W$ through $z$.  Approximate the arch with a margin; the closed path has nonzero winding throughout a sufficiently small open half-neighborhood.  Include the depth of $W$ itself among the side data.

Consecutive shield arcs attach, so their graph depths differ by a bounded amount, uniformly over exterior samples.  This bounds all differences $d_i-d_1$ within a finite integer set.  For the lower bound in \eqref{wall:shield-formula}, realize a shortest full graph route to $w$ as a continuous finite-arc path.  Its base rim has a point outside the comparison box; a route on $W$ itself already has its recorded depth.  Winding forces the route to have a contact with a shield arc on the very same wall vertex.  To check this equality-contact point, suppose there were no such contact.  A switch between route segments cannot occur at a shield intersection, as it would give a same-kind equality.  On each route segment, touching shield arcs on walls of the other kind stay on the closure of one side of its Jordan wall by noncrossing.  Flatten the wall and perturb those arcs slightly toward that side, avoiding the segment's endpoints and switches.  Compactness permits this locally for each segment.  The perturbation changes neither endpoint index, so the winding index cannot change along the route, a contradiction.  After the last passage into the common box, retain the suffix from such an equality contact.  It gives the lower bound by the definition of $\ell_{\mathrm{loc}}$.  The upper bound is immediate by concatenating a shortest route to the shield with a local route.  At contact occurrences on $W$, use the same direct traversal or nearby open-side points.  Finally retain parity, so an unrecorded common shift can be taken even.
\end{proof}

A \emph{caller} is a disk together with the active state of its recursive fan exploration.

\begin{proposition}[Finitely many local callers]\label{wall:finite-local-patterns}
After further buffering, the following data have only finitely many joint possibilities over the compared conditional extensions: canonical callers that have boundary, trunk, or hit-wall operations near $K$ and interact with the exterior buffer; their local open domains and Jordan boundary arcs; their first two half-layer pieces and relevant child groupings; and their half-layer depths relative to the corresponding local shield.  A separate common shift is suppressed for each local depth array.  Whole internal subtrees have common determined geometry.  The finiteness is pathwise, with a random finite cover and random finite lists; it is not a deterministic bound and does not bound the number of preceding remote-only steps.
\end{proposition}

\begin{proof}
In the neighborhoods of Lemma~\ref{wall:shield-depth}, very low relative levels produce no new local trace, except possibly the already recorded side of $W$.  Empty ancestor sweeps covering the patch are irrelevant.  At very high relative levels there are no active cells passing from an inner buffer out of the comparison box.  Indeed choose a sufficiently late layer in one representative with tiny active mesh, and use \eqref{wall:shield-formula} and the bounded relative alternatives to include that layer locally in every case.  Direct child boundaries already reached there cannot be entered by active cells.

Only boundedly many intermediate relative levels remain.  Their local half-layer sets, holes, and sides are known, for each of finitely many alternatives.  Outside grouping of components may differ, but only finitely many pieces cross a fixed positive buffer.  In a representative at any of these levels there are finitely many macroscopic cells, each Jordan.  Their boundaries have finitely many arc portions crossing a positive-width buffer, by uniform continuity of their Jordan parametrizations.  The same holds for open component passages: a traversing component either has an adjacent boundary passage across a smaller buffer or contains an interior ball of a fixed positive size.  Equivalently use a closed-disk Jordan parametrization and uniform local joinability of nearby interior points.  Thus finitely many external groupings suffice.  Smaller wholly internal pieces have no outside grouping, and very deep internal cells are components of the known local depth pattern with determined subsequent hierarchy.

On $W$, use the two side arrays, starting its interior base afresh when $W$ is a child, or both sides when it is central.  Treat complete nearby child interiors similarly.  A covering cell with no boundary or layer trace present locally has no tip or welding operation there.  Taking a finite subcover of $K$, and if useful of an intermediate enlargement, gives finite lists for every caller crossing the outer buffer with an operation near $K$.  Group matching instances and boundary-arc orders across the cover by finite tags.  Except for finitely many larger crossing cases, a child meeting $K$ is in a common internal subtree or is a child with locally known boundary in one such caller; the Lebesgue covering margin handles small children in one member.  Buffers may be chosen from countable bases with strict margins.  This proves the stated quantifiers.
\end{proof}

\paragraph{Coupling local source and target pools.}\label{wall:source-pools}
Consider the first fan of a compared even-layer caller $C$.  Its oriented interior gaps are fixed by adjacent signs and are all traversed.  Its source is length sampled.  When the Liouville field is also common on the patch, shared boundary arcs with the same interior side have identical length measure, by the prescribed-target covariance argument in the proof of Proposition~\ref{cut:branching-fields}, applied to the actual birth laws: canonical positive disks or finite-branch negatives.

Partition buffered boundary portions into finitely many disjoint measurable categories contained in matched arc pieces, plus a category for the remainder.  Refine for overlapping charts and the finite grouping alternatives.  For each matched caller/category slot, sample one common point from normalized local length; in each extension use it with the category's actual length probability, keeping exterior choices private.  Distinct caller/category slots use independent draws, so each extension has the required independent fresh length-sampled sources.  Match whole internal callers directly, including countably many such domains.  Add common dense length pools on matched arcs and private completions elsewhere; for path equality tests their finite-list laws need only be absolutely continuous with respect to ordinary endpoint length sampling conditional on the fields.  This construction is conditional on the geometry/field array: first assign the finite crossing patterns on a compact enlargement, then allocate disjoint slots to their boundary pieces.  Slots of distinct cells in any one extension remain distinct even at touching arcs.  Diffuse-measure partitions handle category edges; zero-length categories need no point.  Thus each crossing case has only finitely many local source choices, while each internal caller has a shared source.  No fresh Gaussian law is asserted after this coupling; almost-sure path facts transfer from the valid individual marginal and endpoint-density laws.

\begin{lemma}[Finite interior markers]\label{wall:markers}
Between any two fixed nested buffer boundaries within common caller geometry, there is a finite collection of directed interior subarcs, or points on them, such that every tested fan-path passage across the buffers traverses a marker strictly within that passage.  The same collection works for all matched configurations.  Directions, orders and grouping of the markers require only finite tags.
\end{lemma}

\begin{proof}
Away from the caller boundary, common crosscuts and macroscopic arc finiteness give the assertion.  It remains to exclude passages that hug the boundary arbitrarily closely.  Such a sequence crossing a middle band would converge along a nondegenerate Jordan boundary interval, with uniformly vanishing depth.  Choose within that interval a short common full fan gap $g$, with distinct endpoints and lying in the matching region; Lemma~\ref{wall:fan} supplies such gaps arbitrarily locally.

Use a Jordan half-collar away from the ends of the limiting traversed interval.  The intermediate value property gives a subpassage crossing its longitudinal sides in the thin strip.  Put both endpoints of $g$ strictly between those sides, leaving boundary subintervals before, between and after them that $g$ does not hit.  Points near these subintervals are alternately on the two sides of the crosscut, with positive margins.  If the path meets the open gap, it must traverse its interior excursion, contradicting arbitrary shallowness.  Otherwise simplicity forces it to visit both endpoints.  One of these is the terminal endpoint of the oriented excursion through $g$, so every fan branch hitting it shares its prefix, including that excursion, again a contradiction.  Longitudinal crossings follow by shrinking a common Jordan collar, or by lifting to its boundary-parameter interval.

Every crossing passage therefore reaches a uniformly positive interior depth somewhere strictly within it.  Only finitely many common arcs enter this interior core.  Uniform continuity and their path order provide a finite marker selection with margins intercepting every such passage.
\end{proof}
\subsection{Local uniqueness of directed fan portions}

\begin{proposition}[Local path uniqueness]\label{wall:fan-local-uniqueness}
Compare ordinary fan paths in locally identical caller geometry, with the same real field $J-c$.  Consider portions from a common directed interior marker to a common later marker $B$, lying entirely, with room, in the agreement region.  Either endpoint may instead be the shared actual source or shared actual local target.  In the initial-marker case retain one binary tag: the true-boundary direction toward the target at the first boundary contact in the portion, if there is one.  With matching tags the portions coincide.  For a common terminal target $y$ this tag is already determined by the boundary grouping and linear order away from the source $x$: at a contact proceed toward $y$ without passing back across $x$.

The assertion holds simultaneously for the countably many marginal samples and ordinary endpoint tests used above.  It requires no new tag at each later boundary contact.
\end{proposition}

\begin{proof}
We first prove a protected domain-restriction statement for a single path, then use it to compare two paths.  Choose a stopping cut $\sigma$ strictly inside a shared incoming marker, or use the actual start.  Its individual conditional field and future path have the usual level-line prescription in the true remaining disk: constant $c$ on unhit true boundary, oriented incoming slit values when present, and insertion values $c-1,c+1$ on a drawn future prefix.  Cuts can be selected from a countable stopping library of ordered small-circle hits after crossings, with a strict shared past segment.  They are taken separately in each path's valid marginal law.

\paragraph{Protected artificial domains and stopping tests.}
Let $T$ be a simply connected artificial domain sharing a protected range and its incoming tip/slit with the true remaining domain.  Choose an artificial target $z_*$ and data $c$ away from short incoming intervals with values $c-1,c+1$; all true data near the protection match, and $z_*$ is outside the incoming intervals.  Stop when the path leaves the protected range, approaches $z_*$, or its true continuation side toward the actual target first becomes incompatible with the side toward $z_*$ in $T$.  We claim that up to these stops it agrees with the determined ordinary level line in $T$, read from the protected field.

Conditional on the true domain and past at $\sigma$, all tests can be drawn from a countable library: intersect with polygonal range-neighborhoods, take the accessible component, choose dense prime-end targets and incoming-data cutoffs in that component, then choose polygonal protected sets with separated collars.  Only the component through the range and its matching true boundary are used; in a disconnected intersection shrink collars to keep positive clearance from other components.  Domain and data agreement have positive room to the artificial changes.

Admissibility is tested from the current prefix, domain and actual-target side at boundary contacts.  At the endpoint of an incoming open excursion, inspect which of the two adjacent true-boundary directions is open toward the target in each domain minus the prefix.  At a contact accumulated by earlier contacts, the true path can proceed only ahead of them; require the artificial target to remain in the component accessible on that side.  Component membership is a Borel past test, using polygonal paths and prime boundary arcs.  First failure is a stopping time in the completed right-continuous filtration, or is approached by increasing admissible stops.  We only need prefixes with room for these tests a little beyond a specified point.

The same tests can be described directly in the initial prime coordinates of $T$.  Use a slightly smaller protection collar, avoid both targets and, in the marker case, the old source.  Incoming intervals cover only the needed portions of the common slit.  Every later hit is on unhit true boundary, never on the old path.  Along such an approach the two domains share an open Jordan half-neighborhood, so their prime-coordinate transition extends to the boundary, by reflection in prime coordinates.  This also holds at the incoming slit and its tip, away from its distant end.  Thus the protected path has a continuous simple lift to $T$.

Stop at the first true-boundary contact where the direction toward the actual target in the original $x,y$ linear order differs from the direction toward $z_*$ in the artificial source-target order.  Away from starts and targets this direction comparison is locally constant on tested contacts.  It agrees near a common true start, and the incoming slit is never revisited.  Hence a first discrepancy with clearance is attained; protection exits may be stopped first.

Before discrepancy, the future really lies in the artificial target component.  At the end of an open excursion, Jordan separation places it next to the short untouched boundary ray in the matched direction.  At an accumulation contact that forward ray is untouched and accessible, and at most one complementary component is accessible there by a continuation avoiding the prefix.  Indeed the earlier excursions are complete disjoint crosscuts with diameters tending to zero; two distinct components sharing the point would need a separating crosscut ending at it, as in the shadow construction in Lemma~\ref{wall:fan}.  No later tested contact can lie behind an earlier one on the same artificial flank: close the path from the artificial source to that contact outside the prime disk to separate the forward ray from the proposed later contact.  The future avoids that path and its endpoints.  Applying this at every earlier contact yields monotone progress toward $z_*$, an untouched target arc and the required accessible component.  At first mismatch this remains true up to the contact itself, after which an artificial continuation can take the artificial side.

\paragraph{Continuous Loewner growth in the protected range.}
We give the elementary boundary check needed for the stopped characterization.  In initial prime-disk coordinates of $T$, the target component at a boundary visit has Jordan boundary consisting of the traced arc between the last contacts on its two flanks, using the source where necessary, and the untouched target arc.  The two arcs are disjoint except at their ends; the path before the earlier of those contacts cannot enter this interior.  At an interior time append the final outgoing simple slit.  Thus the tip has a single access, and short future portions grow through that access.  Continuity of the trace and positive-clearance tests of accessibility from the target give kernel convergence of the mapping-out domains from either time direction.  Capacity is continuous and strictly increasing, since every nontrivial interval contains genuine interior growth; no true boundary interval is traversed at zero capacity.

The driving tip is also continuous at a boundary visit.  In the current target component take a tiny crosscut from a visible point on the incident trace just before the visit to a nearby untouched true-boundary point.  It can shrink both in physical prime coordinates and in the current mapping-out chart.  It cuts off the approaching tips shortly before the visit: the trace since the last opposite-flank visit, the crosscut, and the boundary about the target form a separating curve.  It also cuts off the short future afterward.  At an interior tip use a crosscut between the two incident slit flanks ending strictly earlier on the trace; at the initial tip use the initial boundary.

Their mapped diameters are uniformly small at nearby times.  In bounded-range half-plane coordinates, the Beurling estimate makes the probability that Brownian motion from a high interior point hits a small crosscut before the boundary uniformly small, because that cut is attached to the connected complement.  Standard hull distortion keeps its image bounded.  Conversely a connected arc attached to the real boundary with diameter bounded below has a uniform positive hitting probability: either it spans a fixed scale above the line, using the interior continuum estimate, or it traverses a rectangle of fixed width near the line, which Brownian motion has positive probability to cross vertically.  Thus mapped diameters vanish.  A strict interior point of any fixed cut tracks under kernel convergence, so the cut-off tip preimages squeeze to the same driver from both time directions.  This proves continuous simple Loewner growth before the protected stops, including up to a first incompatible contact.

\paragraph{One positive change of Gaussian domain.}
Apply the two-collar argument of Lemma~\ref{wall:noncontact} to the true remaining domain and $T$.  Here true boundaries may be rough, so we verify the needed energy estimates.  On shared enlargements of a comparison collar project onto their common zero-Dirichlet piece.  The complementary random part is harmonic.  For an orthonormal basis of this harmonic projection, the sum of pointwise squared basis values on an intermediate neighborhood is the Green variance difference.  It is bounded away from the artificial edge, by representing the difference through exit there in bounded-domain coordinates.  Cutoff Caccioppoli estimates, with zero data at the shared true boundary, sum to a finite total Dirichlet norm.  Bounded harmonic differences of boundary means have the same cutoff property: they are exit extensions of bounded data on nonshared edges, zero on the shared true boundary; approximate and apply the same energy estimate.  This proves the Gaussian equivalence through the shared range, including up to rough true boundary.  It also works with a common inserted slit for the comparison between domains, and on separated collars off that slit for comparisons between inserted and base laws.  The Cameron--Martin and Markov arguments consequently give equivalence to a product local class on the two collars.

For a fixed admissible prefix $a$, let $R$ and $R_a$ be respectively the base and inserted $T$-to-true-domain densities on the restriction inside an outer comparison collar.  Both depend only on that outer collar by the Dirichlet Markov property.  In either fixed domain, insertion-to-base densities on the two collars depend only on the inner one.  Product equivalence implies that $R_a/R$ is constant in the collar field.  Positive density versions conditional on the observed prefix obey this identity.  Tilt the joint base field and path once by $R$, and extend outside the comparison region by the conditional $T$ kernel.  That extension is unchanged by protected slit insertion given the inside.  Therefore, at every tested prefix, the extended tilted field has the complete $T$-slit prescription, including cut-off components at a contact; true data themselves are unchanged outside the slit.

The level-line characterization and determination apply in the Dirichlet disk
with finite piecewise constant data, height gap two, and cumulative side
weights at least $-1$ \cite[Theorems~1.1.1--1.1.3 and Proposition~2.1.7]{WangWu}.  The conditional mean martingales and Green decrements are exactly its defining ones.  By the preceding Loewner check, the tilted path equals the $T$ line until protection exit or inadmissibility.  Extra stopping introduces no uniqueness problem: at the stop, continue with the ordinary level rule in the artificial target component, whose field prescription and admissible boundary data preserve those martingales for the concatenation.  This also covers a stop exactly at the first mismatching contact.

The ordinary $T$ line up to protection exit depends only on the protected field and collar.  Indeed its stopped successful-prefix kernel has the unchanged conditional outside field; localize it by Bayes and use path determination, just as in Lemma~\ref{wall:arc-palm}.  Thus the fresh field extension after tilting cannot alter the compared prefix.  Removing the one positive tilt proves the claim as an almost-sure equality test in the original protected field.  The mismatch test itself need not be exterior-measurable.

Versions can be chosen measurably in $T$, its data and the countable test parameters by disk uniformization and conditional sampling.  Conditional on the individual true past, the accessible components of polygonal simply connected range-neighborhoods, with dense targets and incoming cutoffs chosen using that component itself, form the required library.  Identical component and data choices therefore use the same local version in two samples.  Rational ball-entry and radial-hit cuts inside a strict marker avoid any earlier nonshared visit by simplicity.  Exceptional prefixes are excluded in the two separate valid marginal laws, including length-density endpoint laws; no Gaussian law is claimed after conditioning simultaneously on two different true pasts.

\paragraph{The tests rule out divergence.}
Suppose the two portions first diverge at $q$.  Use the shared actual start or choose $\sigma$ within their incoming marker before $q$.  Thicken the common prefix through $q$ to a thin simply connected neighborhood $G$ in the agreement region, avoiding both earlier nonshared pasts and leaving room for short futures of both paths.  Simplicity gives such a Jordan-arc thickening.  Only their common terminal incoming slit extends back from $\sigma$.  The accessible component of the remaining domain intersected with $G$ is a common $T$.  Take $G$ polygonal and simply connected.  Its connected true complement reaches the exterior: for a crossing caller use a common-buffer ball containing neither marked end, since the caller lies on a side of the separating hierarchy; included pasts attach to that complement.  Internal Jordan callers admit smaller such neighborhoods or require no exterior change.  Thus $T$ can be simply connected.  Jordan access and the shared slit give local domain agreement with room, also at $\sigma$ in slit coordinates.  Polygonal and dense prime-end choices make this one of the countable tests for both paths.

Both short futures enter the same component of $T$ cut along the common prefix.  At interior $q$ this is the single open tip neighborhood.  If earlier boundary contacts accumulate at $q$, they approach along one boundary direction away from the actual nonhit target.  Distinct accessible components would require a separating crosscut ending at $q$, whereas there is no final gap then.  At the endpoint of a last gap, the two paths take its same side.  For the first contact after the incoming marker this is exactly the binary tag.  At a later contact the last crosscut has its two true-boundary ends in the shared region.  Opposite sides cannot both lead to the common later marker $B$ by paths staying in that region and avoiding the crosscut: both routes are valid in the common local geometry of one disk and can be pushed inward until $B$, contradicting separation.  For a shared terminal $y$, its boundary order gives the same conclusion.  No other separating prefix excursion contains $q$.  At the true initial source with empty prefix there is a single interior component.  Hence short future interior points $r_1,r_2$, one on each path, lie in one accessible component on the required side.

Join these points to a common open boundary target arc of that component away from the prefix.  Such an arc exists by extending toward the exterior in prime-disk coordinates while avoiding the compact simple prefix, which does not separate the sphere there.  It need not be one of the short incoming-data intervals; those may shrink, and the path never returns to its incoming slit.  Choose $z_*$ with clearance on this arc.  Its use can be stopped early on both futures: the connecting paths from $r_i$ to the target arc avoid $q$ and the prefix with room.  At every preceding contact the chosen germ still reaches the $z_*$ component along the actual protected future to $r_i$ and then that connecting path.  Simplicity and the shared Jordan boundary ensure connectivity off previous slits.  Thus both original paths pass the same admissibility test through and a little beyond $q$.  The common protected local version then makes their prefixes equal there, contradicting first divergence.  This proves the proposition.
\end{proof}

\subsection{All targets, births, and local length data}

\begin{proposition}[Simultaneous target order]\label{wall:simultaneous-targets}
Use the finite caller patterns and coupled pools above.  On a buffered compact observation region, the full family of local target-path portions has only finitely many joint possibilities, including their path order and shared-prefix tree, simultaneously for a dense family of ordinary targets and for length-almost every target.  Matching portions coincide exactly.  No separate finite choice is introduced for each target.  The conclusion holds jointly in a finite cover of comparison charts and is unchanged by increasing or reordering the dense target lists used to represent the fan.
\end{proposition}

\begin{proof}
In one comparison chart take an inner marked region, a larger local-target region, and still larger buffers.  By Lemma~\ref{wall:markers}, choose finitely many oriented anchors intercepting both passages from the inner region out of the local-target region and passages from the closed local-target region through the next buffer.  Take an additional finite outer anchor set when classifying earlier boundary visits.  A simple path uses each anchor at most once, and all paths using an oriented anchor in one fan have the same prefix to it.  Between consecutive anchors, or from a shared local start to an anchor, portions visiting the inner region remain buffered.  Proposition~\ref{wall:fan-local-uniqueness} gives finitely many exact possibilities; source choices are already finite.  Which of the finitely many endpoint slots give buffered or relevant portions is itself a finite tag.  Paths ending outside the local-target region contribute only such anchored portions near the observation.

For an inner anchor $A$, the set of ordinary targets whose paths use it is one open interval $(l_A,r_A)$ in linear boundary order away from $x$, ignoring target-null endpoints.  Preceding contacts bound that untouched interval; target change preserves the prefix within it, and every path through $A$ must share this prefix.  Two separated target intervals would alternate with two prior contacts and be separated by the trace.  We must classify these intervals jointly over all local targets.

Cover the local-target boundary by finitely many buffered arc portions.  The chronological visits there before $A$ have finitely many possible ordered descriptions: insert the additional outer anchors on its unique prefix, and use local uniqueness on the intervening start, anchor and end portions that visit these arcs.  Retain their circular order and the position of $x$, tagging exterior gaps coarsely.  For any target $y_A$ using $A$, visits below $y_A$ in the linear order progress increasingly in time, whereas visits above it progress decreasingly.  For a fixed ordered contact description, two choices of the split cannot move two distinct contacts from one class to the other: those contacts would have to be both increasingly and decreasingly ordered.  Thus the split has only finitely many possibilities, with finite tags for intervening exterior gaps relative to $x$.  Within the local arcs, nearest contacts fix $l_A,r_A$ exactly; an exterior endpoint needs only its finite gap tag.  This classifies the anchor's entire target interval at once.

For a shared ordinary local target $y$, its ordered inner anchors are now known.  The path from its last anchor to $y$ stays in the larger chart, or it would meet another anchor.  It is common by local uniqueness and the tagged boundary order toward $y$.  With no anchor, the whole path stays there.  No tag is added for $y$.  For exterior targets retain the finitely many possible ordered sequences of buffered anchor pieces.  Their entire local family, with exact shared portions and path order, has finitely many joint possibilities.

For several charts choose anchors in each.  External-only transits contribute only a finite ordering and branching grouping of this finite list; local terminal portions are handled in their matching chart.  The countable dense tests and open target-interval agreement give the same assertion for length-almost every ordinary target.  Increasing or reordering the tests changes neither the unranked rooted traversal tree nor these common open-interval prefixes.
\end{proof}

\begin{corollary}[Births and length readouts in matched blocks]\label{wall:birth-length}
The finite local comparison may include positive-loop discoveries, negative cut domains, the positions at which internal children attach to the traversal tree, and inherited target intervals.  For a block retained wholly in the agreement region, it also includes its macroscopic length path, elapsed Cauchy clock, and uncharted future tree with the shared internal sources.  It does not determine the absolute offset from a possibly exterior path origin.  Under a constant local Weyl length scaling, both these lengths and elapsed Cauchy time scale by the same constant.
\end{corollary}

\begin{proof}
Add the positive loops hit along the matched paths.  Loops crossing buffers are macroscopic and form a finite list with finite grouping choices.  Whether one was discovered on an earlier exterior prefix is a finite tag: the record on an anchored prefix is shared by every path through its terminal anchor, while a terminal portion toward a local target remains local.  New contacts on a local segment are visible there.  Small loops meeting the inner observation region have all their first contacts visible on these chronological prefixes or terminal portions.

A negative pocket wholly local on a matched full branch is read as a complementary component of the trunk, previously hit positive loops and true boundary.  Its pinch is the branch's last trunk contact before leaving it.  After a negative cut the branch has no later tip contact with that disk, by the exact loss-arc and no-two-real-times rules: loss interiors do not transmit forward.  Similarly a positive loop is first met at its attachment.  With additional observation buffers, complementary-domain identifications for any finite matched path/order list have only finitely many crossing groupings, by Jordan continuity, macroscopic-piece finiteness and the common local obstacles.  The accessible pinch side follows the direction or is a finite crossing-jump tag.

For an internal return $R$, choose a branch witnessing its birth.  By then its entire boundary is drawn on the trunk and earlier discovered loop arcs.  On the matching branch those obstacles have also been drawn by that time, including crossing loops whose prior discoveries are tagged, and $R$'s interior is again avoided.  Thus $R$ is already disconnected and the pinch still visited on the trunk is its cut point.  It is not a cut for following original targets into its interior, so it is left for a fresh sweep.  A small internal positive child has the corresponding common first contact.  Branches reaching the same child cut have not yet split: until discovery its interior stays in a single component at every earlier split.  These assertions hold simultaneously for tiny children in a covering chart; the finitely many larger crossing children require only the previous finite grouping/order tests.

An intermediate negative carrying original target labels continues from its pinch toward their untouched original open interval, up to a boundary-length null set.  Earlier contacts bound a single such arc.  Non-strict identifications outside untouched open arcs carry no additional interval or positive length by running-record polarity and length nullity.  An empty inherited interval for a complete negative birth is precisely the stabilized return of Proposition~\ref{wall:returns}.

Lengths on shared boundary arcs with known side agree by arc covariance.  Frontier lengths can be recovered by summing inherited arcs: at ordinary pre- and post-states, almost every height on either side lies in a strictly inherited open interval from the original outer boundary or a positive circle, by the scalar stack rule.  Exceptional running-record heights, including real-contact labels, have zero Lebesgue measure.  At jumps use the same pre/post conventions, including the aperture.  Away from finitely many designated crossing jumps, swept changes are the symmetric Cauchy sums on each directed side, by localized absolute continuity.  Elapsed local clock is recovered from small-jump counts, for instance positive jumps first contacted on the segment, as in the intrinsic-clock argument in the proof of Proposition~\ref{cut:critical-cutting}.  Omitting finitely many crossing jumps does not affect that clock limit.  Constant Weyl scaling multiplies length and Cauchy time by the same factor.  Together with the shared internal sources this identifies the stated block and its future tree, while leaving its exterior-origin offset unspecified.  The separate finite-chart and integer-shift bookkeeping will retain that distinction in Section~\ref{sec:kernels}.
\end{proof}
